% =============================================================
%  Chapter — Operational Performance Reporting of HVDC Systems
%  Definitions and calculation formulas per CIGRE TB 956 (2025),
%  which supersedes TB 590. Reference: IEC 60633:2019,
%  IEC TR 62672:2018.
%
%  Note: time values are expressed in decimal hours
%        (e.g. 6 h 30 min = 6.5 h).
% =============================================================

% ---- Outage terms -------------------------------------------
\intermezzo{Performance Reporting}{Availability definitions \& formulas — CIGRE TB 956}

\begin{frame}{Scope of reporting}
  \begin{block}{CIGRE TB 956 (2025) — supersedes TB 590}
    A standardized protocol for reporting the annual operational
    performance (reliability, availability, maintenance) of HVDC systems.
  \end{block}
  \vspace{0.4em}
  \begin{itemize}
    \setlength{\itemsep}{0.4em}
    \item Reporting period is the calendar year (Jan–Dec).
    \item Covers point-to-point links, back-to-back, multi-terminal
          systems and DC grids (LCC and VSC).
    \item Only events causing a loss of transfer capability are counted;
          AC-system outages are recorded but excluded from HVDC
          reliability figures.
    % \item {\color{atGray}All durations are given in \textbf{decimal hours}
    %       (6 h 30 min $= 6.5$ h).}
  \end{itemize}
\end{frame}


% % ---- Symbols ------------------------------------------------
% \begin{frame}{Symbols}
%   \renewcommand{\arraystretch}{1.25}
%   \begin{tabular}{@{}l l@{}}
%     \textbf{Sym.} & \textbf{Description} \\
%     \hline
%     $P_m$   & Rated capacity (MW), excluding redundant equipment \\
%     $P_o$   & Outage capacity — capacity reduction caused by the outage (MW) \\
%     ODF     & Outage derating factor \\
%     AOD     & Actual outage duration (decimal h) \\
%     EOD     & Equivalent outage duration (decimal h) \\
%     PH      & Period hours (calendar hours in reporting period) \\
%     AOH     & Actual outage hours \quad$=\sum$ AOD \\
%     EOH     & Equivalent outage hours \quad$=\sum$ EOD \\
%     EU / EA & Energy unavailability / availability (\%) \\
%     $U$     & Energy utilization (\%) \\
%   \end{tabular}
% \end{frame}




\begin{frame}{Outage classifications}
  \begin{block}{Outage}
    A state in which the HVDC system/station is unavailable at its maximum
    continuous capacity due to an event related to the converter station
    or DC line.
  \end{block}
  \vspace{0.3em}
  \begin{itemize}
    \setlength{\itemsep}{0.5em}
    \item \textbf{Scheduled outage} — planned, or deferrable to a suitable
          time (e.g. preventive maintenance). Reported as ``PM''.
    \item \textbf{Forced outage} — equipment unavailable but not in a scheduled-outage state.
      \begin{itemize}
        \item \emph{Trips} — sudden interruption by automatic protection or manual emergency shutdown.
        \item \emph{Other forced outages} — unexpected problems forcing a capacity reduction without a trip.
      \end{itemize}
    \item \textbf{Forced Outage Rate (FOR)} - Ratio of the number of all complete or partial outages of one converter 
    (or one pole in a bipolar configuration, respectively) tothe whole duration of operation.
  \end{itemize}
\end{frame}

\begin{frame}{Rated, outage capacity \& derating factor}
  \begin{itemize}
    \setlength{\itemsep}{0.5em}
    \item \textbf{Rated capacity $P_m$} — maximum capacity (MW) for continuous operation under normal conditions, excluding redundant equipment.
          If seasonal, the highest value is used.
    \item \textbf{Outage capacity $P_o$} — the capacity reduction (MW) the outage would have caused if the system were operating at $P_m$.
    \item \textbf{Outage derating factor (ODF)} — ratio of outage capacity to rated capacity:
  \end{itemize}
  \vspace{0.4em}
  {\centering\large $\displaystyle \mathrm{ODF} = \frac{P_o}{P_m}$ \par}
  \vspace{0.3em}
  \centering
  {\color{atGray}\footnotesize A full outage gives $\mathrm{ODF}=1$; a partial loss gives $0<\mathrm{ODF}<1$.}
\end{frame}

\begin{frame}{Outage duration}
  \begin{itemize}
    \setlength{\itemsep}{0.6em}
    \item \textbf{Period Hours (PH)}\\Calendar hours in the period: 8760 h/yr (8784 in a leap year)
    \item \textbf{Actual outage duration (AOD)} — time elapsed (decimal h)
          between the start and end of an outage, from the first switching
          action to the return of the equipment to operational readiness.
    \item \textbf{Equivalent outage duration (EOD)} — AOD weighted by the
          derating factor, to account for partial loss of capacity:
  \end{itemize}
  \vspace{0.5em}
  {\centering\large $\displaystyle \mathrm{EOD} = \mathrm{AOD}\times \mathrm{ODF}$ \par}
  \vspace{0.5em}
  \centering
  {\color{atGray}\footnotesize Classified by type as equivalent forced (EFOD)\\ or equivalent scheduled (ESOD) outage duration.}
\end{frame}

\begin{frame}{Time categories\\Actual Outage Hours}
  \begin{itemize}
    \item \textbf{Actual Outage Hours (AOH)} \\
    The sum of actual outage durations within the reporting period

    {\centering $\displaystyle \mathrm{AOH}=\sum \mathrm{AOD}$ \par}

    \item \textbf{Actual forced outage hours (AFOH)} \\
    The sum of equivalent forced outage durations within the reporting period

    {\centering $\displaystyle \mathrm{AFOH}=\sum \mathrm{AFOD}$ \par}

    \item \textbf{Actual scheduled outage hours (ASOH)} \\
    The sum of equivalent scheduled outage durations within the reporting period

    {\centering $\displaystyle \mathrm{ASOH}=\sum \mathrm{ASOD}$ \par}


  \end{itemize}
\end{frame}


\begin{frame}{Time categories\\Equivalent Outage Hours}
  \begin{itemize}
    \item \textbf{Equivalent Outage Hours (EOH)} \\
    The sum of equivalent outage durations within the reporting period

    {\centering $\displaystyle \mathrm{EOH}=\sum \mathrm{EOD}$ \par}

    \item \textbf{Equivalent Forced Outage Hours (EFOH)} \\
    The sum of equivalent forced outage durations within the reporting period

    {\centering $\displaystyle \mathrm{EFOH}=\sum \mathrm{EFOD}$ \par}

    \item \textbf{Equivalent scheduled outage hours (ESOH)} \\
    The sum of equivalent scheduled outage durations within the reporting period

    {\centering $\displaystyle \mathrm{ESOH}=\sum \mathrm{ESOD}$ \par}


  \end{itemize}
\end{frame}


\begin{frame}{Energy unavailability \& availability}
  \begin{itemize}
    \setlength{\itemsep}{0.4em}
    \item \textbf{Energy unavailability (EU)}\\Percentage of energy that could not be transmitted due to outages:
    \vspace{0.3em}

    {\centering\large $\displaystyle \mathrm{EU}\,(\%) = \frac{\mathrm{EOH}}{\mathrm{PH}}\times 100$ \par}
    \vspace{0.3em}

    \item \textbf{Forced Energy Unavailability (FEU)}
    \vspace{0.3em}

      {\centering\large $\displaystyle \mathrm{FEU}=\frac{\mathrm{EFOH}}{\mathrm{PH}}\times 100$ \par}

    \item \textbf{Scheduled Energy Unavailability (SEU)}
    \vspace{0.3em}

    {\centering\large $\displaystyle \mathrm{SEU}=\frac{\mathrm{ESOH}}{\mathrm{PH}}\times 100$ \par}

    \item \textbf{Energy availability (EA)}\\Complement of unavailability:
    \vspace{0.2em}

    {\centering\large $\displaystyle \mathrm{EA}\,(\%) = 100 - \mathrm{EU}$ \par}
  \end{itemize}
\end{frame}

\begin{frame}{Energy utilization}
  \begin{itemize}
    \item \textbf{Energy utilization $U$} — measure of the energy actually
          transmitted over the system, relative to its maximum possible:
  \end{itemize}
  \vspace{0.5em}
  {\centering\large
    $\displaystyle U\,(\%) = \frac{\mathrm{Total\,Energy\,Transmitted}}{P_m \times \mathrm{PH}}\times 100$ \par}
  \vspace{0.6em}
  \begin{itemize}
    \setlength{\itemsep}{0.3em}
    \item Total energy transmitted $=$ energy exported $+$ energy imported
          (MWh), referred to the point at which $P_m$ is defined.
  \end{itemize}
\end{frame}


% --------------------------------------------------
\intermezzo{HVDC projects}{Guarantees and Liquidated damages}

\begin{frame}{Typical values for a HVDC project}
  \centering
  \renewcommand{\arraystretch}{1.4}
  \begin{tabular}{llcc}
    \toprule
    \textbf{Metric} & \textbf{Description} & \textbf{Typical value} & \textbf{Guaranteed} \\
    \midrule
    FEU & Forced Energy Unavailability     & 0.2 - 1.0\,\%      &  0.5\,\%\\
    SEU & Scheduled Energy Unavailability  & 1 - 3\,\%          &  1.0\,\%\\
    FOR & Forced Outage Rate (pole) & $\approx$ 4 - 8 / yr &  2 - 3 / yr\\
        & Forced Outage Rate (bipole)      & $\lesssim$ 0.5 / yr  &  0.1 / yr\\
    \bottomrule
  \end{tabular}

  \vspace{1em}
  {\color{atGray}\footnotesize
    \begin{itemize}
      \item Indicative ranges from CIGRE biennial HVDC reliability surveys.
      \item Actual figures vary with technology (LCC vs.\ VSC), scheme age, and transmission medium — long overhead lines raise forced-outage counts, cable schemes lower them.
      \item For bipole outages values of 0.05 / y have been guaranteed, but the values do no account the real performance and it is assumed that liquidated damages are included in commercial risk.
    \end{itemize}   
  }
\end{frame}

\begin{frame}{Liquidated damages — FOR \& SEU}
  \centering
  \renewcommand{\arraystretch}{1.4}
  \begin{tabular}{llcc}
    \toprule
    \textbf{Metric} & \textbf{Penalty rate}\\
    \midrule
    FOR & 0.1 - 1.2 Mio. \$\\
    FEU & 75 - 600 tsd. \$\\
    \bottomrule
  \end{tabular}

  \vspace{1em}
  \begin{itemize}
    \setlength{\itemsep}{0.3em}
    \item Penalties (liquidated damages) apply when the measured value exceeds the contractually guaranteed threshold.
    \item Typically assessed per unit of excess
    \begin{itemize}
      \item FOR: additional forced outage
      \item FEU: addional 0.1\,\% of unavailability
    \end{itemize}
    \item The cap liquidated damages is in the range of 3-5\,\% of the contract value
  \end{itemize}

\end{frame}





